\documentclass[12pt]{article}
\usepackage[utf8]{inputenc}
\usepackage[T1]{fontenc}
\usepackage{hyperref}
\title{Wogamusin Language Resources for K-12 Curriculum Development}
\author{Compiled by Claude Sonnet 5.5 (claude-sonnet-5-5) for Andr\'as Kornai}
\date{}
\begin{document}
\maketitle

\section{Introduction}
Wogamusin (ISO 639-3: wog; Glottocode: woga1249) is a Sepik language of the Ambunti District, East Sepik Province, Papua New Guinea. Wikipedia reports about 700 speakers in four villages (1998). Joshua Project gives a higher figure of 1,600 for the people group. Sources also differ on its subgrouping within Sepik. It is not known to be taught in schools. Very little exists for this language.

\section{Primary Educational Resources}
None found.

\section{Linguistic and Reference Materials}
\begin{itemize}
\item Wikipedia article on Wogamusin, which notes noun classes reminiscent of Bantu: \url{https://en.wikipedia.org/wiki/Wogamusin_language}
\item Wikipedia article on the Wogamus languages: \url{https://en.wikipedia.org/wiki/Wogamus_languages}
\item Glottolog entry: \url{https://glottolog.org/resource/languoid/id/woga1249}
\item Numeral data page from the Max Planck Institute for Evolutionary Anthropology: \url{https://lingweb.eva.mpg.de/channumerals/Wogamusin.htm}
\end{itemize}
No dictionary or grammar was found.

\section{Religious and Community Resources}
Joshua Project reports that Scripture translation has started, with limited resources. It lists no complete Bible, New Testament or audio Bible. \url{https://joshuaproject.net/languages/wog}

\section{Digital and Audio Resources}
None found.

\section{Curriculum Development Context}
No published teaching or reference materials were found beyond encyclopedic summaries. A curriculum would need new documentation work and primers developed with the community. Speaker numbers should be verified locally, given the conflicting figures.

\end{document}